\clearpage\pagestyle{otMain}\setlength{\otBlocEspacement}{2.5mm}\fontsize{11}{15}\selectfont\corrigesChapitre{7}{Proportions et pourcentages}
\begin{blocLivre}{cor-7-1}\begin{corrige}{1}{Lire une grille}\demarche{Compter les cases sur cent.}
42 cases sur 100 : $\res{\frac{42}{100}}=\res{0{,}42}=\res{42\text{~\%}}$.
\label{fig-cor-7-1}\figureGeo[corrige]{pct-1}
\end{corrige}\end{blocLivre}
\begin{blocLivre}{cor-7-2}\begin{corrige}{2}{Trois écritures}\demarche{Passer par une fraction de dénominateur 100.}
$\frac34=\res{0{,}75}=\res{75\text{~\%}}$ ; $0{,}1=\res{\frac1{10}}=\res{10\text{~\%}}$ ; 60~\% $=\res{\frac35}=\res{0{,}6}$ ; $\frac15=\res{0{,}2}=\res{20\text{~\%}}$.
\end{corrige}\end{blocLivre}
\begin{blocLivre}{cor-7-3}\begin{corrige}{3}{Pourcentages usuels}\demarche{Associer chaque pourcentage à une fraction simple.}
La moitié de 46 : $\res{23}$. Le quart de 80 : $\res{20}$. Le dixième de 350 : $\res{35}$. Trois quarts de 12 : $\res9$.
\end{corrige}\end{blocLivre}
\begin{blocLivre}{cor-7-4}\begin{corrige}{4}{Une proportion}\demarche{Partie sur tout, puis dénominateur 100.}
$\frac9{25}=\frac{9\times4}{25\times4}=\frac{36}{100}$ : $\res{36\text{~\%}}$ des élèves.
\end{corrige}\end{blocLivre}
\begin{blocLivre}{cor-7-5}\begin{corrige}{5}{Partie d’un tout}\demarche{Appliquer le pourcentage au tout.}
$\frac{30}{100}\times200=\res{60}$ billes rouges : 30 billes pour chaque centaine.
\end{corrige}\end{blocLivre}
\begin{blocLivre}{cor-7-6}\begin{corrige}{6}{Des proportions simples}\demarche{Multiplier ou diviser pour obtenir un dénominateur 100.}
$\frac7{20}=\frac{35}{100}$ : $\res{35\text{~\%}}$. $\frac{12}{50}=\frac{24}{100}$ : $\res{24\text{~\%}}$. $\frac34$ : $\res{75\text{~\%}}$. $\frac{45}{300}=\frac{15}{100}$ : $\res{15\text{~\%}}$.
\end{corrige}\end{blocLivre}
\begin{blocLivre}{cor-7-7}\begin{corrige}{7}{Appliquer à une grandeur}\demarche{Garder l’unité de la grandeur.}
$\frac15\times3{,}5=\res{0{,}7\ \mathrm{kg}}$. $\frac{15}{100}\times60=\res{9\ \mathrm{min}}$. 10~\% de 240 vaut 24, donc 5~\% vaut $\res{12}$ €.
\end{corrige}\end{blocLivre}
\begin{blocLivre}{cor-7-8}\begin{corrige}{8}{Les trajets du matin}\demarche{Le tout représente 100 \%.}
À pied : $100-45-30=\res{25\text{~\%}}$. Bus : $\frac{45}{100}\times240=\res{108}$ élèves ; vélo : $\res{72}$ ; à pied : $\res{60}$. Contrôle : $108+72+60=240$.
\label{fig-cor-7-8}\figureGeo[corrige]{pct-8}
\end{corrige}\end{blocLivre}
\begin{blocLivre}{cor-7-9}\begin{corrige}{9}{Comparer des proportions}\demarche{Comparer des pourcentages plutôt que des effectifs.}
A : $\frac{18}{30}=\frac35=\frac{60}{100}$, soit 60~\%. B : $\frac{13}{25}=\frac{52}{100}$, soit 52~\%. La proportion est plus grande dans la $\res{\text{classe A}}$, même si les effectifs diffèrent.
\end{corrige}\end{blocLivre}
\begin{blocLivre}{cor-7-10}\begin{corrige}{10}{Le pourcentage manquant}\demarche{Compléter à 100 \%, puis appliquer.}
Sans avis : $100-46-38=\res{16\text{~\%}}$. Nombre : $\frac{16}{100}\times500=\res{80}$ personnes.
\end{corrige}\end{blocLivre}
\begin{blocLivre}{cor-7-11}\begin{corrige}{11}{Les soldes}\demarche{Calculer la réduction avant de la retrancher.}
Réduction : $\frac{30}{100}\times60=\res{18}$ €. Nouveau prix : $60-18=\res{42}$ €.
\end{corrige}\end{blocLivre}
\begin{blocLivre}{cor-7-12}\begin{corrige}{12}{Le jus de fruits}\demarche{Décomposer 12 \% en 10 \% plus 2 \%.}
10~\% de $1{,}5$ L valent $0{,}15$ L et 2~\% valent $0{,}03$ L : 12~\% valent $\res{0{,}18\ \mathrm{L}}$, soit $\res{18\ \mathrm{cL}}$.
\end{corrige}\end{blocLivre}
\begin{blocLivre}{cor-7-13}\begin{corrige}{13}{La commune et sa forêt}\demarche{Appliquer chaque pourcentage au même tout.}
Forêt : $\frac{35}{100}\times2\,400=\res{840}$ ha. Champs : $\res{960}$ ha. Reste : $100-35-40=25$~\%, soit $\res{600}$ ha. Contrôle : $840+960+600=2\,400$.
\end{corrige}\end{blocLivre}
\begin{blocLivre}{cor-7-14}\begin{corrige}{14}{L’élection des délégués}\demarche{Simplifier la proportion avant de la convertir.}
Inès : $\frac{14}{28}=\frac12$, soit $\res{50\text{~\%}}$. Jules et Sami : $\frac7{28}=\frac14$, soit $\res{25\text{~\%}}$ chacun.
\end{corrige}\end{blocLivre}
\begin{blocLivre}{cor-7-15}\begin{corrige}{15}{Vrai ou faux}\demarche{Écrire chaque calcul avec une fraction.}
20~\% de 50 : $\frac{20}{100}\times50=10$. 50~\% de 20 : la moitié de 20, soit 10. C’est $\res{\text{vrai}}$ : dans les deux cas, on calcule $\frac{20\times50}{100}$.
\end{corrige}\end{blocLivre}
